\chapter{特征工程与特征选择：什么对模型真正重要}

特征工程是把领域认识翻译为模型可利用表示的过程。它既不是单纯增加列数，也不是让 AI 枚举所有可能组合，而是在可获得信息、预测时点和模型归纳偏置之间建立联系。实践教材通常将特征工程与特征选择视为一个迭代过程：构造、评估、诊断和修订相互交替，而不是在建模前一次完成\cite{kuhn2019feature}。

原始数据很少能够直接喂给模型。一张记录了每五分钟流量和速度的检测器表格，在送入任何机器学习算法之前，必须先被转化为模型能够理解和利用的数字表示。这个转化过程，就是特征工程。

特征工程在整个建模流程中承担的角色，比它通常得到的重视要大得多。一个平凡的模型配上精心设计的特征，往往优于一个复杂的模型配上粗糙的输入。原因并不神秘：模型能学到的，只有数据里已经存在的信号。如果一个关键的规律没有被编码进特征，模型无论多复杂都无法凭空发现它。

本章的目标是系统回答三个问题：如何从原始数据中构造有价值的特征，如何判断哪些特征值得保留，以及如何在这一过程中与 AI 有效协作而不被它误导。


% ----------------------------------------------------------------
\section{特征工程的定位与价值}
% ----------------------------------------------------------------

\subsection{连接原始数据与模型的桥梁}

许多模型只能从被显式编码的关系中学习。例如，线性模型不会自然理解 23:00 与 0:00 在时间上相邻；不含空间结构的表格模型也不会自动知道两条相连路段可能共享拥堵传播过程。特征工程通过周期编码、滞后量、邻接统计和领域指标，把这些假设写入数据表示。正因为特征包含假设，生成特征后必须说明其计算公式、信息来源、可用时点和预期作用。

特征工程处于整个建模流程的中间位置：它的输入是经过清洗的原始数据，它的输出是模型可以直接读取的特征矩阵。这个位置决定了特征工程的双重依赖性——它既依赖上游的数据质量，又决定下游模型所能看到的信息边界。

在交通工程的实际建模任务中，这条链条通常如下所示。采集的原始数据可能包括检测器的时间戳和流量计数、GPS 轨迹点的经纬度序列、公交 IC 卡的刷卡记录、事故报告的文本描述，或者网约车订单的起讫点和时刻。这些原始字段本身并不包含“早高峰”“节假日前一天”“距离最近地铁站步行五分钟”这样的信息——这些信息需要被主动构造出来，才能成为模型可以学习的信号。

\subsection{AI 能做什么，人必须做什么}

特征工程处于一个特殊的分工地带。在这里，AI 和人类各自有不可替代的贡献，但职责的边界必须划清楚。

AI 擅长的部分包括：将特征构造想法翻译成可运行的代码，对已知特征类型（时间、距离、滚动统计）提供规范化的实现，帮助检查构造逻辑是否存在语法错误或计算偏差，以及在给定特征列表后生成批量处理脚本。

人必须承担的部分包括：决定哪些领域规律值得被编码为特征，定义时间窗口和阈值的业务含义（“早高峰”到底是 7:00—9:00 还是 7:30—9:30），判断某个特征是否会带入未来信息从而引发数据泄露，以及在统计结果与业务经验相冲突时做出最终裁决。

这个分工意味着，特征工程的核心价值不是代码，而是决策。代码可以让 AI 写，决策必须由人来定。

\subsection{好的特征的三个标准}

在开始构造之前，明确判断标准有助于后续的筛选决策。一个好的特征应当满足以下三个条件。

\textbf{与目标相关}：特征应当与预测目标之间存在真实的信息关联，而不仅仅是统计上的偶然相关。例如，在预测路段拥堵时，“是否为工作日早高峰”比“检测器编号的奇偶性”更可能与拥堵存在因果关联。\textbf{不引入未来信息}：特征必须严格只使用在预测时刻之前已经可以获得的信息。用未来时刻的流量计算“当前时段”的均值，会让模型在训练时获得它在部署时绝不可能拥有的信息，导致评估结果虚假乐观。这个问题将在后续专节中详细讨论。\textbf{对模型可用}：特征需要以模型可以处理的格式提供。大多数经典机器学习模型无法直接处理原始的时间戳字符串、中文分类标签或不定长文本。这些原始信息需要经过数值化或编码之后，才能成为可用特征。

\subsection{特征太少与特征太多的双重风险}

初学者面临的最常见错误有两种，它们方向相反，但都会损害模型性能。

特征太少会导致欠拟合。如果用来预测公交客流的特征只有时间戳，模型无法区分工作日与节假日、晴天与暴雨天、开学季与寒暑假，预测精度将大幅受限。领域中真实存在的影响因素，必须通过特征才能传递给模型。

特征太多则会带来一组不同的问题。无关或冗余特征会引入噪声，干扰模型对真实信号的学习；高维特征空间在样本量有限时会加剧过拟合风险；大量高度相关的特征会影响线性模型的系数估计稳定性；过多特征还会降低模型的可解释性，使得向业务方说明“为什么做出这个预测”变得困难。

合理的特征数量没有统一标准，但一个实用的原则是：先从领域知识出发构造有明确业务含义的特征，再用数据驱动的方法评估和筛选，而不是一开始就把所有能想到的特征全部加入。

\subsection{特征工程是迭代过程}

特征工程不是在建模开始前一次性完成的准备工作。更准确的认识是：它贯穿整个建模周期，在每一轮实验中根据模型反馈不断调整。

一个典型的迭代路径如下：初始版本只包含最基础的时间和空间特征；模型训练后，分析残差分布，发现节假日期间的预测误差明显偏大；据此构造节假日前后的特殊标记特征；下一轮训练后，SHAP 分析显示某个“高峰期”标记的贡献与预期相符，但某个 POI 距离特征的贡献接近零；重新审查发现该 POI 类型与目标场景关联性较弱，将其替换为更合适的变量。这个过程会持续，直到模型性能与可解释性都达到可接受的水平。

这意味着在项目中保留特征构造历史记录非常重要，包括每个特征的构造动机、具体定义、引入时间和后续评估结论。AI 可以帮助生成并维护这份文档，但需要人来持续判断哪些特征值得保留。


% ----------------------------------------------------------------
\section{从原始数据构造特征}
% ----------------------------------------------------------------

\subsection{时间特征的构造}

时间是交通数据中最普遍、信息量最丰富的维度之一。几乎所有交通现象——拥堵、客流、事故率——都表现出显著的时间规律性。将时间戳转化为结构化特征，是交通建模的基础操作之一。

\paragraph{基础时间特征}

最直接的时间特征是从时间戳中提取各个时间维度的数值：小时（0--23）、星期几（0--6）、日期、月份（1--12）、季度（1--4）以及年份。

\begin{quote}
我有一列时间戳字段 \texttt{timestamp}，格式为 \texttt{YYYY-MM-DD HH:MM:SS}。请用 pandas 提取以下时间特征并添加到 DataFrame 中：小时（\texttt{hour}）、星期几（\texttt{weekday}，0 为周一）、月份（\texttt{month}）、季度（\texttt{quarter}）。不要修改原始 \texttt{timestamp} 列。
\end{quote}

然而，直接使用这些数值存在一个结构性问题：对于树模型，数值特征可以正常处理，但对于线性模型和神经网络，数值大小被隐含地赋予了顺序关系。更严重的是，小时 23 与小时 0 在物理上是相邻的（深夜到凌晨），但数值上相差最远。

\paragraph{周期性特征的正弦余弦编码}

解决时间周期性问题的标准方法是正弦余弦编码（sin/cos encoding）。将小时 $h$（取值 0 到 23）编码为两个连续特征：

\[
  \text{hour\_sin} = \sin\!\left(\frac{2\pi \cdot h}{24}\right), \quad
  \text{hour\_cos} = \cos\!\left(\frac{2\pi \cdot h}{24}\right)
\]

这样，小时 0 和小时 23 在编码后的二维空间中距离很近，正确地反映了它们在时间上的连续性。类似地，星期几可以用 7 为周期进行编码，月份用 12 为周期。

\begin{quote}
请为以下时间字段添加正弦余弦周期性编码。字段包括：\texttt{hour}（周期 24）、\texttt{weekday}（周期 7）、\texttt{month}（周期 12）。对每个字段生成对应的 \texttt{\_sin} 和 \texttt{\_cos} 两个新列，保留原始数值字段。请使用 NumPy 实现，并添加注释说明编码公式。
\end{quote}

需要注意：正弦余弦编码对树模型（如随机森林、XGBoost）的提升通常有限，因为树模型的分裂操作本身可以处理小时 23 和小时 0 的相邻关系。但对于线性回归和深度学习模型，这种编码往往有明显的改善效果。

\paragraph{节假日与工作日标记}

中国节假日体系具有特殊性：不仅有法定节假日，还有调休安排，导致部分周六、周日实为工作日，而部分工作日是实际休息日。直接用星期几判断是否为工作日，会在调休日前后产生系统性错误。

正确的做法是维护一份官方发布的节假日日历，并据此构造以下特征：\texttt{is\_holiday}（是否为法定假日，含调休后的休息日）、\texttt{is\_workday}（是否为实际工作日，含调休工作日）、\texttt{holiday\_type}（假日类型，如春节、国庆、清明等）、\texttt{days\_to\_holiday}（距离下一个节假日的天数，用于捕捉节前出行激增效应），以及 \texttt{days\_from\_holiday}（距离上一个节假日结束的天数，用于捕捉节后返程效应）。

在 Python 中，\texttt{chinese\_calendar} 库提供了经过维护的中国节假日数据，可以直接调用。

\begin{quote}
请使用 \texttt{chinese\_calendar} 库，为 DataFrame 中的 \texttt{date} 列（日期格式）添加以下特征：\texttt{is\_holiday}（布尔型）、\texttt{is\_workday}（布尔型）、\texttt{holiday\_name}（节假日名称，非节假日为空字符串）。请处理调休情况，确保调休工作日被标记为工作日而非周末。
\end{quote}

\paragraph{高峰期标记}

早晚高峰是交通领域最具预测价值的时间特征之一，但其定义因城市和道路类型而异。在构造高峰期特征之前，应先通过探索性分析（EDA）确认本地的实际高峰时段，而不是简单套用教科书上的“7:00--9:00 和 17:00--19:00”。

不同场景下的高峰时段差异显著。主干道的早高峰通常在 7:30--9:30，晚高峰在 17:00--19:30；商业区周边的晚高峰可能延至 20:00--21:00；学校周边的早高峰可能提前，并出现放学小高峰（约 11:30 和 16:30）；机场快速路则无典型双峰结构，航班时刻的影响更大。可以构造一个有序的类别特征（早高峰、平峰、晚高峰、夜间）而非简单的二元标记，以保留更多信息给模型。

\paragraph{特殊事件特征}

某些周期性更长或一次性的事件对交通产生显著影响，仅靠常规时间特征无法捕捉。\textbf{春运}通常在春节前 40 天开始，形成持续约两个月的高强度长途出行需求，影响城际道路、高铁和机场周边路段；可以构造“距春运开始天数”和“春运阶段”（前段/节中/后段）特征。\textbf{开学季}发生在每年 9 月初和 3 月初前后，学生流集中出现在大学周边区域，校园周边交叉口短时饱和度显著上升。\textbf{大型活动}（体育赛事、演唱会、展览等）会在场馆周边短时集聚数万人，产生极端峰值流量；如果数据集时间范围内包含此类事件，应添加对应的标记特征。

这类特征的构造需要事先整理事件日历，并将其与数据时间范围匹配。AI 可以帮助实现匹配逻辑和日期范围计算，但哪些事件值得纳入，必须由研究者根据任务场景判断。

\subsection{空间特征的构造}

交通现象在空间上同样呈现强烈的规律性：不同等级道路的拥堵模式不同，靠近商业中心的路段流量特征不同于居住区，地铁站出入口周边在高峰期面临完全不同的压力。将空间信息转化为可用特征，是交通建模的另一个核心环节。

\paragraph{道路属性特征}

直接从道路数据库或 GIS 文件中提取的道路属性，通常是最基础的空间特征，主要包括以下几类：道路等级（高速公路、快速路、主干路、次干路、支路，通常编码为有序整数或独热编码）、限速（法定最高车速，单位 km/h）、车道数（每方向车道数）、路段长度（检测器覆盖的路段长度，影响流量计数的空间代表性），以及道路曲率（弯道比例，与事故率相关）。

\paragraph{距离特征}

距离关键地点的远近，往往对应着不同的出行需求结构和流量特征。常用的距离特征包括：距 CBD（中央商务区）的距离（捕捉通勤方向性和流量等级）、距最近地铁站的距离（影响公交—小汽车换乘行为）、距机场的距离（影响深夜和清晨的流量分布）、距高速出入口的距离（影响货运和长途出行的路网选择），以及距最近医院的距离（在事故严重程度建模中影响救援响应时间）。

这些距离可以是直线距离，也可以是路网距离，后者计算成本更高但通常更贴近实际出行行为。在使用 AI 生成这类特征时，需要明确指出距离类型和坐标系：

\begin{quote}
我有两个 DataFrame：\texttt{roads}（包含路段中点的经纬度 \texttt{lon, lat}）和 \texttt{stations}（包含地铁站的经纬度 \texttt{lon, lat}）。请为每条路段计算到最近地铁站的直线距离（单位：米），使用 Haversine 公式处理球面距离。将结果存为 \texttt{dist\_to\_metro} 列，追加到 \texttt{roads} DataFrame 中。坐标系为 WGS84。
\end{quote}

\paragraph{兴趣点（POI）特征}

路段周边的土地利用和兴趣点分布，反映了该区域的功能属性，进而影响出行需求的类型和时间分布。常见的 POI 特征包括：周边学校数量（500 米范围内，与早晚高峰的接送车流相关）、周边医院数量和床位规模（影响非规律性应急出行）、周边商场建筑面积（影响购物型出行和夜间流量）、周边停车场容量（影响小汽车使用率和拥堵形成时间），以及周边住宅小区户数（影响居住型通勤出行的起始量）。

POI 数据通常来自高德、百度地图 API 或 OpenStreetMap，在使用前需要确认数据时效性，因为建筑环境的变化会导致旧版 POI 特征失效。

\paragraph{路网拓扑特征}

道路在路网中的位置决定了其交通负荷的潜在来源。处于路网中心的节点承担更多的过境流量，而处于路网末梢的节点更多服务于本地出行。常用的拓扑特征有以下几种。\textbf{介数中心性（Betweenness Centrality）}衡量通过该节点的最短路径数量，高介数中心性的路口是路网的关键咽喉点，出现事故或施工时影响范围最广。\textbf{紧密中心性（Closeness Centrality）}是到其他节点的平均最短路径长度的倒数，反映该路段到其他地点的可达性。\textbf{入度/出度}表示交叉口连接的进入和离开方向数，影响信控设计的复杂程度和延误分布。\textbf{路网密度}是单位面积内的路网总长度，反映区域可达性水平。

路网拓扑特征的计算通常借助 \texttt{networkx} 或 \texttt{osmnx} 库。下面是一个向 AI 请求生成计算代码的示例 Prompt：

\begin{quote}
我有一份路网图数据，已经存储为 \texttt{networkx} 的有向图对象 \texttt{G}，节点为路口，边为路段。请计算每个节点的介数中心性（betweenness centrality）和紧密中心性（closeness centrality），并将结果存储到 DataFrame \texttt{node\_features} 中，列名分别为 \texttt{betweenness} 和 \texttt{closeness}。如果图较大（节点数超过 5000），请使用近似算法（\texttt{k=500} 的采样版本），并在注释中说明这个选择对结果精度的影响。
\end{quote}

\subsection{滞后特征与滚动统计}

时序交通数据具有强烈的时间自相关性：当前时刻的路段状态，高度依赖于过去若干时刻的状态。滞后特征（Lag Features）和滚动统计（Rolling Statistics）是捕捉这种时序依赖性的标准方法。

\paragraph{为什么需要滞后特征}

如果不引入任何历史信息，模型只能依赖静态属性（时间标记、空间属性）来预测当前的交通状态。但交通拥堵具有明显的传播性和惯性：上一个五分钟的排队长度，往往是下一个五分钟排队长度的最强预测因子之一。将历史状态直接作为输入特征，能够让模型以最直接的方式利用这种时序规律。

\paragraph{常见的滞后时间窗口}

滞后特征按时间跨度可分为几类。短期滞后（5 分钟前、10 分钟前、15 分钟前）用于捕捉拥堵传播和消散的即时动态；中期滞后（30 分钟前、1 小时前）用于捕捉高峰形成与消退的较慢过程；同期历史（昨天同时段、上周同星期同时段）用于捕捉日内周期性和周内周期性；季节性历史（去年同期）则适用于需要捕捉年度周期性的长期预测任务。

\paragraph{滚动统计特征}

在滞后值的基础上，可以进一步计算过去若干时段的统计量。滚动均值（如过去 30 分钟的平均速度）反映当前时段的整体状态水平；滚动标准差（如过去 30 分钟速度的波动程度）在数值较高时可能预示着交通状态的不稳定；滚动最大值和最小值用于捕捉最近时段内的极端状态；滚动变化量（当前值与过去若干时段均值的差）则反映当前状态与近期平均的偏离程度。

\paragraph{关键警告：滞后特征与数据泄露}

\begin{quote}
\textbf{这是特征工程中最容易犯的错误，也是后果最严重的错误之一。}

构造滞后特征时，必须确保所有滞后值都基于预测目标时刻之前已经可以观测到的数据。如果在滚动窗口计算中包含了目标时刻或之后时刻的数据，就发生了数据泄露，模型的训练误差会虚假地降低，但在真实部署时性能会急剧下降。

\textbf{错误做法：}在构造第 $t$ 时刻的特征时，用 $t-2$ 到 $t+2$ 的窗口计算均值，并将其作为预测 $t+1$ 时刻的特征之一。这个均值包含了 $t+1$ 时刻之前的信息，在实际预测时不可获得。

\textbf{正确做法：}在构造第 $t$ 时刻的特征时，只使用严格早于 $t$ 的历史数据。计算滞后值时使用 \texttt{shift(n)} 操作（$n \geq 1$），计算滚动统计时使用 \texttt{rolling(window).mean().shift(1)} 确保窗口不包含当前时刻。
\end{quote}

下面是一个明确声明防泄露约束的 Prompt 示例：

\begin{quote}
我需要为时序交通数据构造滞后特征，目标是预测 \texttt{t+1} 时刻的流量。数据以 5 分钟为粒度，按检测器编号分组。

请为 \texttt{flow} 字段构造以下特征，严格确保不发生数据泄露，即所有特征只能使用严格早于当前时刻的历史数据：
\begin{itemize}
  \item \texttt{flow\_lag1}：前 1 个时段（5 分钟前）的流量
  \item \texttt{flow\_lag3}：前 3 个时段（15 分钟前）的流量
  \item \texttt{flow\_lag12}：前 12 个时段（1 小时前）的流量
  \item \texttt{flow\_roll\_mean6}：过去 6 个时段（30 分钟）的滚动均值
  \item \texttt{flow\_roll\_std6}：过去 6 个时段的滚动标准差
\end{itemize}

对每个检测器独立计算。窗口的计算必须使用 \texttt{shift(1)} 确保不包含当前时刻。请在代码中添加注释说明如何防止数据泄露，并说明为什么 \texttt{rolling().mean()} 之后必须加 \texttt{shift(1)}。
\end{quote}

\subsection{交互特征}

交互特征（Interaction Features）将两个或多个原始特征组合成一个新特征，用于捕捉这些因素联合作用时的效果。

\paragraph{何时应当构造交互特征}

盲目构造所有可能的特征组合会产生海量冗余特征，引入大量噪声，并使模型解释变得困难。构造交互特征的合理前提，是有明确的领域假设支撑某种联合效应的存在。

在交通工程中，以下几类交互具有清晰的领域解释。\textbf{降雨天气 × 早晚高峰}：雨天对拥堵的加剧效应在高峰期更为显著，因为高峰期路网本身已处于饱和边缘；而夜间低流量条件下，雨天的影响相对有限。\textbf{节假日 × 长途路段}：节假日出行以长途为主，高速公路和城市快速路的流量特征在节假日与平日的差异，远大于城市支路。\textbf{学校周边 × 放学时间}：16:00--17:00 的学校周边路段，与其他时间或其他区域的路段，具有根本不同的流量构成（大量接送车辆）。\textbf{赛事活动 × 场馆周边}：大型体育赛事结束后，场馆周边路段在短时内面临极端的散场流量，这种效应高度局部化，交互特征比两个独立特征更能精确描述它。

构造方式通常是将两个特征相乘（对于连续变量和连续变量），或将类别变量的独热编码与连续变量相乘（捕捉分组效应）。


% ----------------------------------------------------------------
\section{领域知识是特征工程的核心}
% ----------------------------------------------------------------

\subsection{为什么 AI 无法自动生成“好”的领域特征}

一个常见的误解是：既然 AI 能够处理各种格式的数据，它是否可以自动从原始数据中发现所有有价值的特征？

答案是否定的，原因在于特征工程本质上需要“将领域知识翻译成计算逻辑”，而领域知识中的大部分并不存在于 AI 的训练数据中。AI 知道正弦余弦编码是处理时间周期性的通用方法，但它无法知道：在你的城市，“节假日前一天”的出行需求是呈现激增还是提前消化，因为这依赖于当地的出行文化和政策环境；AI 知道距离特征是空间建模的常见操作，但它无法知道：在你研究的路网中，距离“快速公交（BRT）站”的步行距离，比距离“普通公交站”的距离对客流的解释力更强，因为这依赖于你所在城市的公交网络结构。

\textbf{领域知识的作用，是将“可能有用的变量”的搜索空间从无穷大压缩到有限的、有业务意义的候选集。}没有这种压缩，特征工程就变成了在黑暗中摸索。

\subsection{如何把领域假设翻译成 Prompt}

有效的人机协作路径是：人提供领域假设（“我认为某种特征可能有价值，因为……”），AI 将其翻译成具体的计算实现。下面通过两个例子说明这个翻译过程。

\paragraph{示例一：节假日前后出行需求不对称}

领域假设：在中国的交通出行中，“节假日前一天”（人们急于赶在假期前回家或出发）和“节假日后第一个工作日”（集中返程后进入工作状态）具有显著不同于普通日期的出行特征，且两者的需求结构也不对称——节前以长途出行为主，节后以城内短途出行为主。

从这个假设出发，可以构造以下特征：\texttt{is\_pre\_holiday}（是否为节假日前一天）、\texttt{is\_post\_holiday}（是否为节假日后第一个工作日）、\texttt{days\_to\_next\_holiday}（距下一个节假日的天数，作为连续变量捕捉“节前效应”的渐进增强），以及 \texttt{days\_from\_last\_holiday}（距上一个节假日结束的天数，捕捉“节后恢复”的衰减过程）。

将这个假设告诉 AI，让它实现具体计算：

\begin{quote}
我的 DataFrame 有一列 \texttt{date}（日期格式）和一列 \texttt{is\_holiday}（布尔型，已标记法定节假日和调休）。

根据以下领域假设构造特征：中国交通出行中，节假日前一天和节后第一个工作日具有特殊的出行特征，且随距离节假日天数呈现渐变效应。

请构造：（1）\texttt{is\_pre\_holiday}，为 True 当且仅当次日为节假日第一天；（2）\texttt{is\_post\_holiday}，为 True 当且仅当昨日为节假日最后一天；（3）\texttt{days\_to\_next\_holiday}，到下一个节假日第一天的天数（当天为节假日则为 0）；（4）\texttt{days\_from\_last\_holiday}，距上一个节假日最后一天的天数（当天为节假日则为 0）。使用 pandas 实现，按时间顺序处理，不要使用循环遍历每一行。
\end{quote}

\paragraph{示例二：公交站点可达性影响客流}

领域假设：公交站点周边的客流量，不仅取决于时间因素，还取决于步行可达范围内的居住和就业人口密度。一个位于居住区内部、步行 5 分钟覆盖大量住宅的站点，其早高峰客流特征，与一个位于商业区边缘、步行覆盖范围内只有办公楼的站点截然不同。

据此，可以为每个公交站点构造以下特征：\texttt{pop\_500m}（站点 500 米步行圈内的居住人口，来自人口普查数据）、\texttt{emp\_500m}（站点 500 米步行圈内的就业岗位数量，来自企业注册数据或土地利用数据）、\texttt{res\_building\_area\_500m}（站点 500 米范围内住宅建筑面积总量），以及 \texttt{commercial\_area\_500m}（500 米范围内商业用地面积）。

AI 可以帮助实现“以站点为圆心、给定半径的空间连接”这一计算逻辑，但“500 米是合理的步行可达范围阈值”这一判断，来自领域知识而非数据本身。

\subsection{特征创意的来源}

好的特征创意并非凭空产生，以下几种来源值得系统利用。\textbf{文献回顾}：同类研究（特别是发表在交通工程顶级期刊上的实证研究）中使用的变量列表，是特征候选集的重要来源；如果已有研究反复证明某类变量与目标高度相关，忽视它需要充分的理由。\textbf{EDA 发现}：在探索性数据分析中，如果某两个变量的散点图呈现出规律性的非线性关系，或者某个子组在目标变量上表现出显著的分布差异，这往往预示着相应的交互特征或分组特征具有价值。\textbf{业务经验}：长期从事路段管理、公交调度或交通执法的从业人员，往往能够凭直觉指出哪些因素在实践中确实重要；将这些经验系统地转化为特征，是研究者与实践者合作的核心价值。\textbf{利益相关方访谈}：在实际项目中，向交通管理部门、出租车司机、公交公司运营人员等利益相关方了解他们认为影响某种交通现象的关键因素，往往能够获得书面文献中未曾系统记录的现场知识。


% ----------------------------------------------------------------
\section{特征选择：不是所有特征都有价值}
% ----------------------------------------------------------------

特征构造完成后，面对的下一个问题是：哪些特征应当保留，哪些应当剔除？特征选择不是特征工程的附属步骤，而是其必要组成部分。

\subsection{为什么要做特征选择}

保留所有构造出来的特征，通常不是最优策略。无关特征会引入噪声，干扰模型的信号提取，在训练样本有限时尤为突出；高度相关的冗余特征会影响线性模型的系数估计稳定性，导致系数解释失效；过多特征增加计算成本，拖慢训练速度，在需要频繁重训练的生产系统中尤为明显；冗余特征还会稀释特征重要性的统计显著性，使得 SHAP 值等可解释性工具的输出更难解读；在向非技术受众解释模型决策时，特征数量越少越容易说清楚“模型依据了哪些因素”。

\subsection{特征选择的三类方法}

特征选择研究通常把方法划分为过滤法、包裹法和嵌入法，并强调选择策略应同时考虑泛化性能、计算成本和变量解释\cite{guyon2003introduction}。任何重要性排序都依赖具体数据、模型和评价过程；它回答的是“在当前设定下哪些变量帮助模型”，而不是“现实世界中哪些因素最重要”。

\paragraph{过滤法（Filter Methods）}

过滤法独立于模型，直接根据特征与目标之间的统计关系评分，然后按分数筛选。常用的过滤法指标有以下几种。\textbf{Pearson 相关系数}适用于连续特征与连续目标，衡量线性相关程度，但它无法捕捉非线性关系，相关系数为零并不意味着特征无用。\textbf{Spearman 秩相关}对非正态分布和单调非线性关系更鲁棒，适合交通数据中常见的有偏分布。\textbf{互信息（Mutual Information）}衡量特征与目标之间的广义信息相关性，能够捕捉非线性关系，对分类和回归任务均适用。\textbf{方差阈值法}则直接删除方差接近零的特征，因为这类特征几乎不包含信息（如“是否在研究时段内”这种几乎全为 True 的标记）。

过滤法速度快，但忽略了特征之间的交互效应：两个单独与目标弱相关的特征，组合起来可能具有强预测力。

\paragraph{包裹法（Wrapper Methods）}

包裹法将特征子集的选择视为一个搜索问题，通过训练模型、评估性能来评价不同特征组合的优劣。

最常用的包裹法是递归特征消除（Recursive Feature Elimination, RFE）：训练一个模型，根据特征重要性排名，删除最不重要的特征，在缩减后的特征集上重新训练，重复这个过程，直到剩余特征数量达到目标值。

包裹法比过滤法能捕捉到更复杂的特征交互，但计算成本高，对于大特征集可能不实际。

\paragraph{嵌入法（Embedded Methods）}

嵌入法在模型训练过程中同时完成特征选择，将选择逻辑内嵌于学习算法本身。\textbf{LASSO 正则化}在线性回归的损失函数中加入 $L_1$ 正则项，会将不重要特征的系数压缩到精确为零，从而自动实现稀疏特征选择；通过调节正则化强度参数 $\lambda$，可以控制被保留特征的数量。\textbf{树模型特征重要性}：随机森林、梯度提升树（XGBoost、LightGBM）等模型在训练过程中记录每个特征在所有树的分裂中降低杂质度（Gini 不纯度或均方误差）的总贡献，据此生成特征重要性排名。\textbf{基于 SHAP 的选择}将 SHAP 值的绝对均值作为特征重要性度量，比传统的树模型内置重要性更稳定，且能处理特征之间的相互作用。

\subsection{SHAP 值的获取与解读}

SHAP 将合作博弈中的 Shapley value 用于解释单次预测，使不同特征归因方法能够在统一框架下比较\cite{lundberg2017unified}。但 SHAP 解释的是模型如何使用特征，而不是特征对现实结果的因果作用；当特征高度相关时，归因还会受背景分布和实现选择影响。因此，SHAP 图适合用于发现模型行为和生成检查问题，不应被直接写成政策结论。

SHAP（SHapley Additive exPlanations）是目前最广泛使用的模型可解释性工具之一。它将模型的每次预测分解为每个特征的贡献量，正值表示该特征使预测值升高，负值表示使其降低。

\begin{quote}
我已经训练好了一个 XGBoost 模型 \texttt{model}，训练数据为 \texttt{X\_train}（DataFrame，列名为特征名）。请使用 \texttt{shap} 库完成以下操作：（1）计算 SHAP 值；（2）绘制特征重要性摘要图（\texttt{summary\_plot}），按绝对 SHAP 值均值从大到小排列，颜色表示特征值的高低；（3）绘制 \texttt{hour\_sin} 和 \texttt{flow\_lag1} 两个特征的依赖图（\texttt{dependence\_plot}），显示 SHAP 值随特征值的变化关系。所有图表保存为 PNG 文件，dpi=150。
\end{quote}

在解读 SHAP 摘要图时，应注意以下几点：SHAP 值高度依赖于模型本身，不同模型对同一特征给出的 SHAP 值可能有所不同；高 SHAP 重要性意味着该特征对模型预测影响大，但不直接等于该特征与目标的因果关系；当两个特征高度相关时，SHAP 值会在它们之间分配，导致单个特征的重要性被低估；依赖图（dependence plot）中如果出现非线性、非单调的关系，说明模型学到了复杂的特征效应，值得结合领域知识深入分析。

\subsection{统计重要性与业务必要性的冲突}

特征选择的最终决策不能只依赖统计结果。统计重要性和业务必要性是两个不同维度，两者的冲突需要由研究者做出判断。

\textbf{统计上“不重要”但业务上不可缺少的特征}在以下情形中出现：某个特征在当前数据集上恰好方差较小（例如研究时段内几乎没有雨天），因此统计重要性低；但在实际部署时，模型将在各种天气条件下使用，因此“降雨量”特征必须保留。另一种情况是，某个特征对于满足监管要求或审计需要是必要的，无论其统计表现如何。

\textbf{统计上“重要”但可能是噪声代理变量的特征}也值得警惕。如果一个特征与目标变量统计相关，但这种相关性只存在于特定历史时段而不具有泛化性，那么保留它可能导致部署后性能大幅下降。例如，某个特殊事件导致数据集中某段时间的交通特征异常，与之巧合相关的特征可能获得虚假的高重要性。

\subsection{特征选择决策的记录与说明}

特征选择的结果应当被文档化，包括：最终保留了哪些特征，为什么保留；哪些候选特征被剔除，剔除的理由是什么（统计原因还是业务判断）；不同特征组合的实验结果对比。这份记录既是科学严谨性的要求，也是向合作者和审稿人说明建模决策合理性的依据。


% ----------------------------------------------------------------
\section{防止数据泄露：特征工程中最危险的错误}
% ----------------------------------------------------------------

\subsection{什么是数据泄露}

数据泄露的本质，是训练或评估过程使用了模型真实部署时不可获得的信息，因而得到过度乐观的性能估计。系统研究表明，泄露既可能来自特征本身，也可能来自数据划分、重复样本、时间关系和预处理流程\cite{kaufman2012leakage}。它之所以危险，是因为代码通常能够正常运行，指标甚至异常优秀；只有回到预测时点和信息可得性，才能发现逻辑错误。

数据泄露（Data Leakage）是指在构造特征或训练模型的过程中，无意间将目标变量的信息，或只有在预测时刻之后才能获得的信息，混入了模型的输入中。

数据泄露最险恶之处在于：它不会在训练阶段产生任何警告，反而会让训练误差和验证误差都异常漂亮，直到模型部署到真实环境中，才会发现预测性能远低于评估时的表现。从结果倒推原因，往往要花费大量时间。

\subsection{交通预测中的典型数据泄露场景}

\paragraph{用包含目标时刻的窗口计算“历史”统计量}

假设目标是预测 $t+1$ 时刻的流量，构造了“过去一小时的平均流量”特征，但在计算时，窗口的截止时刻是 $t+1$ 而不是 $t$，即使用了 $t-11$ 到 $t+1$ 共 12 个时段的均值。这个均值中包含了目标时刻本身的流量值，相当于模型在训练时获得了它在部署时绝不可能拥有的“答案的一部分”。

在 pandas 中，这种错误通常表现为：使用了 \texttt{rolling(12).mean()} 而没有追加 \texttt{.shift(1)}。

\paragraph{用测试集的统计量对训练集进行标准化}

另一类常见泄露发生在数据预处理阶段。如果先对整个数据集（包括测试集）计算均值和标准差，再进行 Z-score 标准化，测试集的分布信息就已经被编码进了训练集的特征中。正确做法是：只在训练集上拟合 \texttt{StandardScaler}，然后用同一个 scaler 变换测试集，绝不能用测试集更新 scaler 的参数。

\paragraph{用目标变量本身或其强代理变量构造特征}

如果预测目标是路段拥堵指数（TCI），而构造的某个特征是“当前时段检测器报告的排队长度”，而排队长度实际上几乎等于 TCI 的另一种表达，这会导致模型性能虚假地高，但这个特征在部署时（拥堵指数是预测目标而非已知量时）无法使用。

\paragraph{跨检测器的泄露}

在空间聚合特征的构造中，如果目标是预测某条路段的状态，而特征中包含了相邻路段在“同一时刻”（而非“当前时刻之前”）的状态，且这两条路段之间有强烈的实时同步关系，则相邻路段的同期状态相当于部分泄露了目标路段的状态。

\subsection{如何在 Prompt 中声明防泄露约束}

每次请求 AI 构造时序特征时，应明确将防泄露要求写入 Prompt：

\begin{quote}
以下所有特征的构造，必须严格保证不发生数据泄露。具体规则：（1）所有滞后特征和滚动统计，只能使用严格早于当前行时间戳的历史数据，即 \texttt{shift(n)}（$n \geq 1$）必须在滚动计算之后使用；（2）标准化所用的均值和标准差，只能在训练集上计算；（3）空间聚合特征中，相邻路段的状态只能使用其上一时刻的值，而非当前时刻。请在代码注释中为每个关键操作标注“防泄露”说明。
\end{quote}

\subsection{检验是否发生数据泄露的方法}

以下几种方法可以帮助发现潜在的数据泄露：

\begin{enumerate}
  \item \textbf{权衡检验：}如果模型在验证集上的表现异常好（如 $R^2 > 0.99$ 或均方误差极低），而数据质量并不出众，应首先怀疑是否发生了泄露。
  \item \textbf{特征与目标的相关系数检验：}对所有特征与目标变量计算相关系数。如果某个特征与目标的相关系数过高（如 $|r| > 0.95$），应仔细审查它是否是目标的代理变量或包含了目标信息。
  \item \textbf{时间打乱测试：}将测试集的时间顺序打乱后重新评估，如果性能没有显著下降，说明时序特征可能没有真正利用时序信息，或者时序信息来自泄露而非真实的预测能力。
  \item \textbf{逐特征消融：}将特征一个一个移除，重新评估性能。如果移除某个特征后性能急剧下降，而这个特征在业务逻辑上不应对预测有如此巨大的贡献，应仔细检查其构造逻辑。
\end{enumerate}


% ----------------------------------------------------------------
\section{用 Prompt 管理特征工程流程}
% ----------------------------------------------------------------

\subsection{分步 Prompt 策略}

面对一个完整的特征工程任务，不应该用一个超长的 Prompt 要求 AI 一次性生成所有特征的代码。更可靠的方式是分阶段推进：先规划特征清单，再逐组实现，最后验证输出逻辑。

\textbf{第一步：生成特征规划文档}

\begin{quote}
我正在构建一个预测城市路段日均拥堵时长的回归模型。数据来自某市 2022 年全年的固定检测器数据，字段包括：时间戳（5 分钟粒度）、检测器编号、速度（km/h）、流量（辆/5 分钟）、占有率（\%）。

请按照时间特征、空间特征、滞后特征、统计特征、交互特征五个类别，列出你建议构造的特征清单，每个特征包括：（1）特征名称；（2）构造方法（一句话）；（3）预期与目标的关联方向（正/负/非线性）；（4）数据泄露风险等级（低/中/高）及说明。不要写代码，只输出这份规划表格。
\end{quote}

拿到特征规划表之后，由研究者根据领域知识对其进行增删和调整，确认后再进入实现阶段。

\textbf{第二步：分组实现特征代码}

每次只请求 AI 实现一个类别的特征，并在 Prompt 中明确引用上一步确认的规划内容：

\begin{quote}
根据上面确认的特征规划，请实现“滞后特征”类别中的全部特征。数据已存储在 DataFrame \texttt{df} 中，按 \texttt{detector\_id} 和 \texttt{timestamp} 排序。函数接受 DataFrame 作为输入，返回添加了所有滞后特征的 DataFrame。请遵守防泄露规则（所有滚动统计后必须 shift(1)），并在函数开头添加文档字符串说明每个特征的含义。
\end{quote}

\textbf{第三步：验证特征逻辑}

\begin{quote}
请为上面生成的特征构造函数编写验证代码。验证内容包括：（1）对一个已知的示例行，手动计算各特征的预期值，并与函数输出进行比较；（2）检查任何特征列是否在 \texttt{t=0}（数据集开始处）存在不应出现的非 NaN 值；（3）打印前三行的所有新特征值，供人工检查。如果发现不一致，请报告具体的不一致内容而不是静默通过。
\end{quote}

\subsection{让 AI 输出特征文档}

特征文档是长期项目中的重要资产，记录了每个特征的含义、构造方法、引入版本和历史评估结果。AI 可以根据已实现的代码自动生成初版特征文档：

\begin{quote}
请阅读以下特征构造函数的代码，为其中每个新增的特征列生成一条文档条目，格式如下：

特征名 | 数据类型 | 取值范围 | 构造方法简述 | 预期业务含义 | 注意事项

注意事项应特别指出该特征是否依赖外部数据（如节假日日历）、是否存在数据泄露风险、以及在数据集开头若干行是否会产生 NaN 值。
\end{quote}

\subsection{特征版本管理}

在迭代建模过程中，不同实验往往使用不同的特征集。建立简单的特征版本管理机制，能够避免“这次实验用的是哪个版本的特征”这种令人沮丧的困惑。

一个轻量级的管理方法是：将每次实验的特征集名称（如 \texttt{feature\_set\_v3}）、包含的特征列表和对应的评估结果，记录在一个统一的实验日志文件（CSV 或 JSON）中。AI 可以帮助生成维护这份日志的脚本，以及在不同特征版本之间进行对比的分析代码。

\begin{quote}
请生成一段 Python 代码，将以下实验结果记录到日志文件 \texttt{experiment\_log.csv} 中：特征集版本名称（字符串）、包含的特征名列表（逗号分隔）、验证集 RMSE、验证集 MAE、训练时间（秒）、备注（字符串）。如果日志文件不存在则创建并写入表头，如果已存在则追加新记录。时间戳自动从系统时间获取。
\end{quote}


% ----------------------------------------------------------------
\section{实战练习}
% ----------------------------------------------------------------

\subsection*{练习一：时间特征构造与周期性编码}

某市有一份道路检测器数据集，字段包括：\texttt{timestamp}（格式 \texttt{YYYY-MM-DD HH:MM:SS}，5 分钟粒度）、\texttt{detector\_id}、\texttt{speed}（km/h）、\texttt{flow}（辆/5 分钟）。

\begin{enumerate}
  \item 从 \texttt{timestamp} 中提取以下基础时间特征：小时、星期几、月份、季度。
  \item 对小时（周期 24）和星期几（周期 7）分别构造正弦余弦编码。用 Prompt 请求 AI 实现，然后验证：小时 0 和小时 23 的编码距离，是否确实小于小时 0 和小时 12 的编码距离？
  \item 使用 \texttt{chinese\_calendar} 库，为数据集添加 \texttt{is\_holiday}、\texttt{is\_workday} 字段，并处理调休情况。在你所用的数据集时间范围内，有哪些日期属于“调休工作日”（原本是周末但调休为工作日）？
  \item 绘制“平均速度 vs 小时”的折线图，分别对工作日和节假日绘制两条曲线，观察正弦余弦编码是否合理地反映了两者之间在深夜—凌晨时段的连续性。
\end{enumerate}

\subsection*{练习二：滞后特征构造与数据泄露检验}

使用练习一的数据集，任务是预测每个检测器下一个 5 分钟（$t+1$）的流量。

\begin{enumerate}
  \item 为 \texttt{flow} 字段构造以下滞后特征：\texttt{flow\_lag1}（$t-1$）、\texttt{flow\_lag3}（$t-3$）、\texttt{flow\_lag12}（$t-12$，即 1 小时前）以及过去 6 个时段的滚动均值 \texttt{flow\_roll\_mean6}。
  \item 故意构造一个包含数据泄露的版本：计算滚动均值时不加 \texttt{.shift(1)}，即窗口包含当前时刻。将两个版本的数据集分别训练一个线性回归模型，比较训练集和测试集上的均方误差。泄露版本的性能提升幅度有多大？
  \item 计算所有特征与目标变量 \texttt{flow\_next}（即 $t+1$ 时刻的流量）的 Spearman 相关系数。泄露版本中，哪个特征的相关系数异常高？这是否与你的预期一致？
  \item 写一段用于防泄露自检的验证代码：对数据集最开头的若干行，手动验证 \texttt{flow\_lag1} 的值是否确实等于上一行的 \texttt{flow}，以及 \texttt{flow\_roll\_mean6} 的值是否确实是前 6 行（不含当前行）的均值。
\end{enumerate}

\subsection*{练习三：特征重要性分析与业务判断}

选取一份交通事故数据集（可使用公开的 NGSIM、UK Road Safety、或城市交通管理部门开放的事故报告数据），目标是预测事故的严重程度（轻伤/重伤/死亡，多分类）。

\begin{enumerate}
  \item 从原始数据中构造至少 15 个特征，涵盖时间（小时、星期、节假日）、空间（道路等级、弯直路段、限速）、天气（降雨、能见度、路面状况）和交通状态（速度、流量、饱和度）四个维度。写出每个特征的构造动机，而不仅仅是代码。
  \item 训练一个随机森林分类器，使用内置的特征重要性排名列出前 10 个特征。再使用 SHAP 值重新排名，比较两种方法的结果是否一致。如果存在排名差异，分析可能的原因。
  \item 假设利益相关方（交通执法部门）要求在模型输出中保留“道路限速”特征，因为它在法规执行中具有标准化意义。但 SHAP 分析显示该特征的重要性处于末尾 20\%。你会保留还是删除这个特征？写出你的判断理由，并说明你会如何在报告中向技术受众和非技术受众分别说明这个决策。
  \item 将最终选定的特征集写成特征文档表格（参考本章介绍的格式），包括：特征名、类型、取值范围、构造方法、预期与目标的关联方向、注意事项。
\end{enumerate}

\subsection*{练习四：端到端特征工程 Prompt 链设计}

不使用任何代码，只设计一套完整的 Prompt 链，用于完成以下任务：预测城市主干道上下班高峰期（7:00--9:30 和 17:00--20:00）的平均拥堵指数（连续变量，回归任务），数据为某市 2023 年全年的 5 分钟级检测器数据。

\begin{enumerate}
  \item 设计一个“特征规划 Prompt”，要求 AI 列出候选特征清单（包括时间、空间、滞后和交互特征），并对每个特征评估业务合理性和数据泄露风险。
  \item 设计一个“特征实现 Prompt”，为时间特征类别写出完整的约束条件，包括：周期性编码、节假日处理、防泄露声明和代码规范要求。
  \item 设计一个“特征验证 Prompt”，要求 AI 生成能够发现数据泄露的自检代码，并说明每项检查针对的是哪类泄露风险。
  \item 设计一个“特征文档 Prompt”，要求 AI 根据已实现的代码自动生成特征文档，并指定文档格式和必须包含的字段。
  \item 审视你设计的这四个 Prompt：哪些部分依赖了你的领域知识（交通工程背景），哪些部分是可以泛化到其他领域的通用要求？将它们分别标注出来，并思考：如果把这套 Prompt 链用于航空延误预测，哪些部分需要修改，哪些可以直接复用？
\end{enumerate}





